\thispagestyle{empty}
\pdfbookmark[0]{Ausgabe und Quellen}{translation-sources}
\begin{center}\textbf{Ausgabe und Quellen}\end{center}
\noindent Diese KI-Übersetzung ins heutige Deutsch bietet die astronomischen Schriften aus Carl Friedrich Gauß’ \textit{Werken}, Band~VI (Göttingen, 1874): Abhandlungen, Beobachtungen, Briefauszüge und Besprechungen. Die historischen Titelblätter und das Porträt begleiten den Text, ebenso Ernst Scherings eigene Bemerkungen. Die Übersetzung folgt dem Text der Werkausgabe mit ihren historischen Auszügen und Anmerkungen. Bei deutschsprachigen Originalen handelt es sich um eine Wiedergabe im heutigen Deutsch. Die bandinternen Verweise behalten die Seitenzahlen der Ausgabe von 1874. Quellen und Parallelübersetzungen sind in den begleitenden Anmerkungen angegeben. Die Titelvignette stammt aus dem Digitalisat der \href{https://archive.org/details/werkecarlf06gausrich/page/n8/mode/1up}{University of California Libraries}.
\par\medskip\noindent\textbf{Elliptische Elemente der Pallas}
\par\smallskip
\noindent Diese KI-Übersetzung umfasst alle fünfzehn Artikel und die Tabellen von Gauß’ \textit{Disquisitio de elementis ellipticis Palladis}. Sie folgt dem lateinischen Text in den \textit{Werken}, Band~VI (Göttingen, 1874), einschließlich des historischen Titels der Abhandlung. Auch die in modernem Deutsch wiedergegebenen Bandtitel und das Porträt stammen aus dieser Ausgabe. Die Abhandlung wurde am 25.~November 1810 vorgelegt und 1811 in Band~I der \textit{Commentationes societatis regiae scientiarum Gottingensis recentiores} gedruckt. \href{https://gdz.sub.uni-goettingen.de/id/PPN236007467}{Göttinger Digitalisat}.
\par\smallskip\noindent Historische Paralleltexte sind der deutsche \href{https://www.digitale-sammlungen.de/de/view/bsb10538613?page=463}{Zeitschriftenauszug vom November 1811}, die von Anton Börsch und Paul Simon herausgegebene deutsche Teilübersetzung der Artikel~10--15 (Berlin, 1887, S.~118--128) und Hale F. Trotters \href{https://www.york.ac.uk/depts/maths/histstat/gausspallas.htm}{englischer Auszug} nach der französischen Übersetzung von Joseph Bertrand im Princeton Technical Report Nr.~5 (1957). Die Ausgabe von Börsch und Simon enthält ein Vorwort von Helmert. Die deutsche Terminologie zur Bahnbestimmung wurde auch mit E. F. M. Jochims \textit{Satellitenbewegung}, Band~III, Stand 9.~Februar 2022, S.~60 und 63, abgeglichen.
\par\smallskip\noindent Übersetzung und Satz wurden mit KI erstellt.
\clearpage
